\section{Seguridad y Mitigación de Riesgos}

La seguridad operativa constituye el eje transversal del proyecto PDA STARS. Se implementó una matriz formal de gestión de riesgos basada en las normas de la Tripoli Rocketry Association (TRA) y la National Association of Rocketry (NAR), complementada con procedimientos de verificación obligatoria en rampa coordinados por el Oficial de Seguridad (RSO).

\subsection{Matriz de Identificación, Evaluación y Mitigación de Riesgos}
\begin{table}[H]
\centering
\scriptsize
\begin{tabularx}{\linewidth}{p{2.8cm}p{2.8cm}ccX}
\toprule
\rowcolor{navyiejaga!15}
\textbf{Riesgo Identificado} & \textbf{Causa Potencial} & \textbf{Nivel Inicial} & \textbf{Nivel Residual} & \textbf{Medida de Mitigación Implementada} \\
\midrule
\textbf{Ignición accidental del motor en rampa} & Descarga electrostática o conexión errónea del ignitor & Crítico & Bajo & Uso obligatorio de llave de seguridad en la estación de ignición a 15 m. Inserción del ignitor únicamente en rampa por el RSO tras despejar el área. \\
\rowcolor{grisclaro}
\textbf{Falla en el despliegue del paracaídas} & Quemadura por gases de eyección o atascamiento en el tubo & Crítico & Bajo & Uso de manta térmica Nomex ignífuga de 15x15 cm, verificación de holgura suave de la ojiva ($0,2\text{ mm}$) y pruebas previas de eyección neumática. \\
\textbf{Inestabilidad aerodinámica en vuelo} & $CG$ retrasado respecto al $CP$ o flexión de aletas & Crítico & Bajo & Validación estricta de margen estático $\ge 2,00\text{ calibres}$ en OpenRocket, recorte de cartón y Swing Test. Aletas fabricadas en PETG-CF rígido. \\
\rowcolor{grisclaro}
\textbf{Deriva excesiva fuera de zona ($R > 250\text{ m}$)} & Viento fuerte en altura o paracaídas sobredimensionado & Alto & Bajo & Calibración del paracaídas a 40 cm ($v_{\text{term}} = 8,00\text{ m/s}$) para aterrizar dentro de los 40-50 s (47,86 s simulados). Monitoreo de anemómetro en tierra antes del despegue. \\
\textbf{Desprendimiento de carga útil en shock} & Falla del anclaje mecánico en la apertura & Alto & Bajo & Ojal de retención macizo integrado en la carcasa PETG probado a 250 N, atado permanentemente al cordón principal de Kevlar de 3,6 m. \\
\rowcolor{grisclaro}
\textbf{Apagado de aviónica por vibración} & Falso contacto en la batería por shock de aceleración & Medio & Bajo & Soporte de batería NiMH soldado y asegurado mecánicamente, circuito en PCB CNC y amortiguación interna con espuma EVA. \\
\bottomrule
\end{tabularx}
\caption{Matriz de identificación, evaluación y control de riesgos operacionales de PDA STARS.}
\label{tab:matriz_riesgos}
\end{table}

\subsection{Procedimientos Operativos de Seguridad (Checklist de Rampa)}
El Oficial de Seguridad (RSO) debe ejecutar y registrar en la bitácora la siguiente secuencia de verificación obligatoria antes de autorizar cualquier lanzamiento:
\begin{enumerate}[leftmargin=1.5em]
    \item \textbf{Inspección Previa en Mesa de Armado:}
    \begin{itemize}
        \item Verificar integridad del fuselaje Kraft BT-80 y ausencia de fisuras en las uniones epóxicas de las tres aletas.
        \item Comprobar holgura de acople de la ojiva (debe deslizar suavemente por gravedad invertida con ligera fricción controlada).
        \item Verificar el trenzado de las 8 líneas de suspensión del paracaídas y anclaje firme del cordón de Kevlar a la bahía de carga útil.
        \item Enrollar la manta Nomex cubriendo el paracaídas e insertarlo sin torsión en el fuselaje.
    \end{itemize}
    \item \textbf{Verificación de Aviónica y Carga Útil:}
    \begin{itemize}
        \item Energizar la carga útil Seeeduino XIAO, constatar los destellos LED indicadores de calibración y cierre del archivo microSD.
        \item Verificar voltaje de batería ($\ge 3,6\text{ V}$) y asegurar la fijación de la bahía con los tornillos de retención.
    \end{itemize}
    \item \textbf{Montaje y Aseguramiento en Rampa:}
    \begin{itemize}
        \item Deslizar el cohete sobre la varilla de lanzamiento de 1,5 m ($\varnothing 6,5\text{ mm}$), comprobando recorrido libre y verticalidad.
        \item Insertar el motor Estes F15-4 en el tubo portamotor y asegurar el retén roscado inferior.
        \item Asegurar que la llave de seguridad del controlador de ignición permanezca en posesión exclusiva del RSO.
    \end{itemize}
    \item \textbf{Protocolo de Ignición y Conteo:}
    \begin{itemize}
        \item Verificar despeje de la zona en un radio mínimo de 15 metros alrededor de la rampa.
        \item Medir velocidad del viento con anemómetro de tierra (suspender si el viento sostenido excede los 5,0 m/s).
        \item Inserción del ignitor comercial por el RSO y conexión de pinzas cocodrilo con polaridad correcta.
        \item Inserción de la llave de seguridad en la estación de ignición, activación de señal acústica de alerta y conteo regresivo audible de 5 segundos.
    \end{itemize}
\end{enumerate}
